%! Logarithmic Function Context and Data Modeling — Unit 5, Lesson 5.7 — Exit Ticket (Answer Key)
\documentclass[10pt]{article}
\usepackage{precalculus-article}
\usepackage{precalculus-key}   % pulls in -boxes; do NOT also load -boxes
\newcommand{\TallMath}[1]{$\displaystyle #1\rule[-1.4em]{0pt}{3.2em}$}

\begin{document}

\pageheader{Unit 5, Lesson 5.7}{Exit Ticket --- Answer Key}
\namedateperiod

\vspace{0.15in}

\begin{enumerate}[leftmargin=1.8em, itemsep=18pt]

\item The table below shows a data set.

\vspace{6pt}
\begin{center}
\renewcommand{\arraystretch}{1.8}
\begin{tabular}{|c|c|c|c|c|}
\hline
\rowcolor{sky}
$x$    & $2$ & $4$ & $8$ & $16$ \\
\hline
$f(x)$ & $1$ & $2$ & $3$ & $4$ \\
\hline
\end{tabular}
\end{center}

\begin{enumerate}[leftmargin=1.5em, itemsep=8pt, label=\alph*.]
  \item What is the proportion factor?

  \vspace{4pt}
  Proportion: \ans{2} \qquad Base $b = $ \ans{2}

  \item Write the logarithmic model.

  \vspace{4pt}
  $f(x) = $ \ans{$\log_2(x)$}

  \item Predict $f(32)$.

  \vspace{4pt}
  $f(32) = $ \ans{$\log_2(32) = 5$}
\end{enumerate}

\item Let $f(x) = 3\ln(x)$.

\begin{enumerate}[leftmargin=1.5em, itemsep=8pt, label=\alph*.]
  \item Evaluate $f(e^2)$.

  \vspace{4pt}
  $f(e^2) = 3\ln(e^2) = 3 \cdot $ \ans{2} $= $ \ans{6}

  \item Evaluate $f(1)$.

  \vspace{4pt}
  $f(1) = 3\ln(1) = 3 \cdot $ \ans{0} $= $ \ans{0}

  \item Solve $f(x) = 6$ for $x$.

  \vspace{4pt}
  $3\ln(x) = 6 \;\Rightarrow\; \ln(x) = $ \ans{2} $\;\Rightarrow\; x = $ \ans{$e^2$}
\end{enumerate}

\end{enumerate}

\end{document}
